\section{Publications}

        \begin{twocolentry} {2026}

\textbf{Double Triangle Annotation: A Scalable Human-in-the-Loop Framework for High-Precision Historical Document Annotation} [\href{https://arxiv.org/abs/2605.25781}{Prépublication}] [\href{https://github.com/nmrenyi/double-triangle-annotation}{Annotation}] [\href{https://github.com/nmrenyi/rosenwald-benchmark}{Benchmark}]

\end{twocolentry}

        \vspace{0.10 cm}
        \begin{onecolentry}
            \textbf{Yi Ren} (prépublication arXiv, 2026)
        \end{onecolentry}
        \vspace{0.10 cm}

        \begin{onecolentry}
        \begin{highlights}
            \item Conception d'un cadre d'annotation autour d'un principe de consensus, appliqué récursivement sur deux niveaux---entre deux MLLMs indépendants (Claude, Qwen, Llama, Grok), puis entre deux systèmes human-in-the-loop construits sur ces modèles---n'escaladant que les désaccords vers un jury humain
            \item Word Error Rate de 0.003 atteint sur les Guides Rosenwald, avec >85\% des champs auto-acceptés par consensus inter-modèles---réduisant la charge totale de relecture humaine de plus de 50\% par rapport à une annotation manuelle complète par un seul annotateur
            \item Publication du premier benchmark d'extraction structurée pour les Guides Rosenwald (13 595 champs annotés sur 60 colonnes) pour soutenir les futurs travaux sur le traitement de documents historiques
        \end{highlights}
        \end{onecolentry}
